\documentclass[conference]{IEEEtran}

\usepackage{amsmath,amssymb}
\usepackage{booktabs}
\usepackage{cite}
\usepackage{graphicx}
\usepackage{microtype}
\usepackage{url}
\usepackage{xcolor}

\graphicspath{{./}{results/plots/}}

% Include the repository figure when present and leave a labeled placeholder
% when this manuscript is compiled separately from the results directory.
\newcommand{\resultfigure}[2]{%
  \IfFileExists{results/plots/#1}{%
    \includegraphics[width=\linewidth]{results/plots/#1}%
  }{%
    \fbox{\parbox[c][0.22\textheight][c]{0.94\linewidth}{\centering
      Figure placeholder\\[2mm]\texttt{results/plots/#1}\\[2mm]#2}}%
  }%
}

\title{Expressibility--Trainability Tradeoffs in Noisy Parameterized
Quantum Circuits: A Crossed Feature-Map and Ansatz Study}

% Replace these placeholders before submission.
\author{\IEEEauthorblockN{Author Name(s)}
\IEEEauthorblockA{Department or Laboratory\\
Institution\\
City, Country\\
email@example.com}}

\begin{document}

\maketitle

\begin{abstract}
Selecting a quantum feature map, trainable ansatz, circuit depth, and
entanglement topology jointly determines both the state-space coverage and
the optimization difficulty of a quantum machine-learning model.  We study
this interaction with a crossed design comprising three feature maps, two
ansatz families, three circuit widths, four depths, three connectivity
patterns, calibrated noisy and noiseless simulation, and ten random seeds.
Expressibility is measured by the Kullback--Leibler divergence between
sampled state fidelities and the Haar distribution; entangling capability by
the Meyer--Wallach measure; and trainability by the variance of a
parameter-shift gradient of a data-encoded, composed feature-map--ansatz
circuit.  Across the available noiseless configurations, mean gradient
variance at ten qubits decreases from 0.123 at one repetition to 0.0496 at
eight repetitions.  The decrease is concentrated in the entangling
\texttt{real\_amplitudes} ansatz: for linear ten-qubit circuits, its three
feature-map combinations lose 92.5--99.3\% of gradient variance by depth
eight, whereas the nonentangling \texttt{rotation\_only} ansatz remains
substantially more trainable.  Hardware-calibrated noise produces an even
sharper topology effect.  At ten qubits, the median retained gradient
variance for linear circuits changes from 0.916 at depth one to 0.545 at
depth eight, while most circular and fully connected circuits collapse
toward zero after routing to the sparse device patch.  These results support
joint, hardware-aware selection of encoding, ansatz, depth, and topology,
but do not support a universal feature-map-independent saturation depth.
\end{abstract}

\begin{IEEEkeywords}
quantum machine learning, parameterized quantum circuits, expressibility,
trainability, barren plateaus, noise-aware compilation
\end{IEEEkeywords}

\section{Introduction}

Parameterized quantum circuits are central to near-term variational
algorithms, but their design presents a tension.  Greater state-space
coverage can improve representational capacity, while sufficiently random
or expressive circuits can concentrate objective values and suppress
gradients~\cite{mcclean2018barren,holmes2022connecting}.  At the same time,
near-term processors impose gate errors and sparse connectivity, making
nominally compact all-to-all circuits costly after routing
\cite{preskill2018nisq}.  Circuit selection must therefore consider
expressibility, trainability, and hardware compatibility together.

Quantum feature maps embed classical samples in Hilbert space and underlie
both quantum-kernel and variational-classifier constructions
\cite{havlicek2019supervised}.  Prior expressibility studies introduced
fidelity-based and entanglement-based descriptors and showed that both can
saturate with circuit depth~\cite{sim2019expressibility}.  The present study
extends this circuit-level perspective in three practical directions.  It
(i) crosses feature maps and ans\"atze instead of pairing them one-to-one,
(ii) measures ansatz trainability after composing it with a feature map
bound to a real data sample, and (iii) evaluates the same design under a
fixed, hardware-calibrated noise model and qubit patch.

The main contributions are as follows:
\begin{itemize}
  \item A factorial comparison of three encoding circuits and two
  structurally distinct ans\"atze over width, depth, topology, and noise.
  \item A reproducible measurement protocol linking feature-map
  expressibility to the gradient variance of the composed classifier
  circuit.
  \item Evidence that depth-induced loss of trainability depends strongly on
  ansatz structure, and that routing topology can dominate this loss under
  calibrated noise.
\end{itemize}

\section{Metrics and Problem Formulation}

\subsection{Expressibility}

Following Sim \emph{et al.}~\cite{sim2019expressibility}, two independent
parameter vectors are sampled for a feature map $U(\boldsymbol{\theta})$,
and the fidelity $F=|\langle\psi(\boldsymbol{\theta})|
\psi(\boldsymbol{\phi})\rangle|^2$ is recorded.  For a Hilbert-space
dimension $N=2^n$, the Haar fidelity density is
\begin{equation}
  P_{\mathrm{Haar}}(F)=(N-1)(1-F)^{N-2}.
\end{equation}
Expressibility is the histogram estimate
\begin{equation}
  E_{\mathrm{KL}}=D_{\mathrm{KL}}
  \left(P_{\mathrm{circuit}}\parallel P_{\mathrm{Haar}}\right).
\end{equation}
A smaller $E_{\mathrm{KL}}$ denotes a more Haar-like, and hence more
expressible, circuit ensemble.

\subsection{Entangling Capability}

For each random parameter draw, the Meyer--Wallach global-entanglement
measure~\cite{meyer2002global} is computed as
\begin{equation}
  Q=2\left(1-\frac{1}{n}\sum_{k=1}^{n}\mathrm{Tr}(\rho_k^2)\right),
\end{equation}
where $\rho_k$ is the one-qubit reduced state.  The reported capability is
the sample mean of $Q$, ranging from zero for product states to one for
maximally entangled single-qubit marginals.

\subsection{Trainability}

The trainability proxy is the variance, over random ansatz initializations,
of one component of the gradient of the local observable $Z_0$.  The first
trainable parameter is differentiated by the parameter-shift rule
\cite{schuld2019evaluating}:
\begin{equation}
  g(\boldsymbol{\theta})=\frac{f(\boldsymbol{\theta}+\frac{\pi}{2}\mathbf{e}_1)
  -f(\boldsymbol{\theta}-\frac{\pi}{2}\mathbf{e}_1)}{2},
  \qquad T=\mathrm{Var}_{\boldsymbol{\theta}}[g].
\end{equation}
Larger $T$ indicates gradients that survive random initialization; values
near zero indicate an effectively flat landscape.  Unlike the other two
metrics, $T$ is evaluated on the feature map followed by the ansatz.  The
feature-map parameters are bound to a training sample, leaving only the
ansatz parameters free.

For comparing noise conditions, retained trainability is
\begin{equation}
  R_T=\frac{T_{\mathrm{noisy}}}{T_{\mathrm{noiseless}}}.
\end{equation}
A value near one indicates little noise-induced loss; a value near zero
indicates that noise has erased parameter sensitivity.

\section{Experimental Methodology}

\subsection{Data Preparation}

The experiment uses the 569-example, 30-feature Breast Cancer Wisconsin
(Diagnostic) binary-classification dataset~\cite{uciWDBC}.  A stratified
80/20 split gives 455 training and 114 test examples.  The implementation
standardizes the raw training features, fits principal-component analysis
to the selected qubit count, standardizes the retained components again,
and multiplies them by $\pi/2$ to obtain encoding angles.  Every transform
is fit only on the training split.  Although a test set is prepared, the
current experiment reports circuit descriptors rather than classifier test
accuracy.

\subsection{Crossed Circuit Design}

The feature-map families are:
\begin{itemize}
  \item \texttt{local\_z}: repeated single-qubit $Z$ rotations without
  entanglement;
  \item \texttt{zz}: a diagonal, commuting $ZZ$ feature map; and
  \item \texttt{cx\_ry}: data-dependent $R_y$ rotations followed by CX
  entanglers.
\end{itemize}
The same data parameters are reused at every repetition.  Each feature map
is crossed with (i) \texttt{rotation\_only}, containing repeated $R_y$
layers and no entanglers, and (ii) \texttt{real\_amplitudes}, containing
$R_y$ layers and CX entanglers.  The crossed design separates the encoding
effect from the ansatz effect more clearly than fixed feature-map--ansatz
pairing.

The planned grid contains widths $n\in\{4,8,10\}$, depths
$d\in\{1,2,4,8\}$, and linear, circular, and full entanglement patterns.
The topology is immaterial for the entirely local
\texttt{local\_z}+\texttt{rotation\_only} circuit, so that pair is evaluated
only with the linear label.  This gives 192 planned circuit configurations.
One result shard---\texttt{local\_z}+\texttt{real\_amplitudes}, 10 qubits,
depth 8, full topology---is absent from the archived results.  The analysis
therefore uses 191 complete configurations, each observed for ten seeds and
two noise conditions, totaling 3,820 raw rows.

\subsection{Simulation and Noise}

Noiseless metrics use exact statevector simulation.  Per seed,
expressibility, entangling capability, and gradient variance use 500, 200,
and 200 Monte Carlo samples, respectively.  Noisy evaluation uses density
matrices and reduced counts of 60, 30, and 30 samples to control runtime.
Ten seeds both initialize the Monte Carlo sampling and select the first ten
training examples cyclically; consequently, variation across seeds includes
both random initialization and data-instance variation.

The implementation uses Qiskit~\cite{javadiabhari2024qiskit} and a noise
model derived from the 127-qubit \texttt{FakeSherbrooke} backend.  Every
circuit is restricted to a nested prefix of physical qubits 0--9, which
forms a connected chain in the selected heavy-hex patch.  Parameterized
circuits are transpiled once with optimization level 3 to the
\{\texttt{sx}, \texttt{rz}, \texttt{x}, \texttt{ecr}\} basis and then
simulated with the backend-derived noise model.  This fixed patch avoids
confounding circuit width with a change of device while exposing routing
overhead when circular or full logical connectivity is embedded in sparse
hardware.

\subsection{Analysis Protocol}

Each reported cell is the mean over ten seeds; plotted error bars are one
sample standard deviation.  Depth trends are descriptive rather than
inferential.  Rank correlations use the 12 linear, noiseless cells formed by
three feature maps and four depths within each ansatz and width.  Missing
configurations are not imputed.

\section{Experimental Results}

\subsection{Feature-Map Expressibility Is Family Specific}

Fig.~\ref{fig:saturation} shows that increasing depth does not produce one
universal expressibility curve.  At ten qubits with linear connectivity,
\texttt{cx\_ry} improves from $E_{\mathrm{KL}}=0.536$ at depth one to
$2.73\times10^{-3}$ at depth eight.  The \texttt{zz} family is already
Haar-like relative to the other families and remains within 0.045--0.105.
By contrast, repeated \texttt{local\_z} rotations do not expand the reachable
product-state family; their estimated divergence increases from 0.364 to
5.79 at depth two before decreasing to 2.81 at depth eight.  These results
show family-dependent plateaus and nonmonotonic finite-sample behavior, not
a common saturation depth.

\begin{figure*}[t]
  \centering
  \resultfigure{poster_saturation.png}{Expressibility and gradient variance
  versus depth at ten qubits with linear topology.}
  \caption{Depth dependence in the noiseless, ten-qubit linear setting.
  Lower KL divergence means greater expressibility; larger gradient variance
  means better trainability.  Solid and dashed trainability curves denote
  \texttt{rotation\_only} and \texttt{real\_amplitudes}, respectively.}
  \label{fig:saturation}
\end{figure*}

\subsection{Ansatz Structure Controls Depth-Induced Gradient Loss}

Table~\ref{tab:endpoints} isolates the ansatz effect at ten qubits with
linear connectivity.  Every \texttt{real\_amplitudes} combination loses
more than 92\% of its noiseless gradient variance between depths one and
eight.  The strongest collapse occurs for \texttt{cx\_ry}, from 0.264 to
$1.98\times10^{-3}$ (99.3\%).  The nonentangling
\texttt{rotation\_only} ansatz is more resistant: its \texttt{cx\_ry}
combination loses 64.0\%, while its \texttt{local\_z} and \texttt{zz}
combinations finish above their depth-one values, albeit nonmonotonically.

\begin{table}[t]
  \caption{Noiseless Gradient-Variance Endpoints (10 Qubits, Linear)}
  \label{tab:endpoints}
  \centering
  \small
  \begin{tabular}{llrrr}
    \toprule
    Feature map & Ansatz & $T_{d=1}$ & $T_{d=8}$ & Change \\
    \midrule
    \texttt{cx\_ry} & real amp. & 0.2637 & 0.0020 & $-99.3\%$ \\
    \texttt{local\_z} & real amp. & 0.1518 & 0.0113 & $-92.6\%$ \\
    \texttt{zz} & real amp. & 0.0987 & 0.0072 & $-92.7\%$ \\
    \texttt{cx\_ry} & rotation & 0.2870 & 0.1032 & $-64.0\%$ \\
    \texttt{local\_z} & rotation & 0.2373 & 0.3275 & $+38.0\%$ \\
    \texttt{zz} & rotation & 0.0988 & 0.1509 & $+52.7\%$ \\
    \bottomrule
  \end{tabular}
\end{table}

After averaging over all available feature-map, ansatz, and topology cells,
ten-qubit noiseless gradient variance decreases from 0.123 at depth one to
0.0496 at depth eight, a 59.8\% reduction.  The corresponding reduction at
four qubits is 36.7\% (0.139 to 0.0879), indicating a width--depth
interaction.  In the linear setting, the descriptive Spearman association
between $E_{\mathrm{KL}}$ and $T$ is positive, as expected when greater
expressibility (smaller KL) accompanies smaller gradients.  Across widths
4, 8, and 10, the correlation is 0.273--0.462 for
\texttt{real\_amplitudes} and 0.860--0.895 for
\texttt{rotation\_only}.  Fig.~\ref{fig:tradeoff} visualizes these
trajectories; their differing shapes caution against reducing trainability
to expressibility alone.

\begin{figure*}[t]
  \centering
  \resultfigure{poster_tradeoff.png}{Noiseless gradient variance versus
  expressibility, faceted by ansatz and qubit count.}
  \caption{Noiseless expressibility--trainability trajectories for linear
  topology.  Rows are ans\"atze, columns are widths, marker shape identifies
  the feature map, and marker size increases with depth.  Both axes are
  logarithmic.}
  \label{fig:tradeoff}
\end{figure*}

\subsection{Sparse Hardware Strongly Penalizes Dense Logical Topology}

Fig.~\ref{fig:noise} reports $R_T$ at ten qubits.  For the six linear
feature-map--ansatz combinations, the median retention is 0.916, 0.873,
0.824, and 0.545 at depths 1, 2, 4, and 8, respectively.  At depth eight,
the three \texttt{rotation\_only} linear circuits retain 0.593--0.930,
whereas the three \texttt{real\_amplitudes} circuits retain 0.134--0.498.

The circular and full logical topologies are qualitatively different after
routing to the chain-like patch.  At ten qubits and depth eight, all three
circular \texttt{real\_amplitudes} combinations have
$R_T<6\times10^{-11}$; the available full-topology combinations are also
effectively zero.  Even for \texttt{rotation\_only}, entanglers in the
feature map can induce this routing penalty: depth-eight circular
\texttt{cx\_ry} retains $7.91\times10^{-4}$ and circular \texttt{zz} is
numerically indistinguishable from zero.  The result is therefore not simply
``noise increases with depth.''  It is an interaction between the composed
circuit's entanglers and the mismatch between logical and physical
connectivity.

\begin{figure*}[t]
  \centering
  \resultfigure{poster_noise_topology.png}{Retained gradient variance under
  calibrated noise versus depth and topology at ten qubits.}
  \caption{Retained trainability $R_T$ under the calibrated Sherbrooke noise
  model.  Color and marker identify the feature map; line style identifies
  logical topology.  The dotted horizontal reference is $R_T=1$.}
  \label{fig:noise}
\end{figure*}

\subsection{Implications for Circuit Selection}

The combined results favor a staged configuration policy.  First reject
topologies that require extensive routing on the target coupling graph.
Then compare feature-map expressibility within the remaining shallow
circuits, and finally evaluate trainability on the actual composed
feature-map--ansatz model.  Feature-map descriptors alone are insufficient:
the same encoding can sit above a robust rotation-only ansatz or an
entangling ansatz whose gradients collapse with depth.  Conversely, a
nominally nonentangling ansatz does not eliminate noise sensitivity if its
feature map introduces dense two-qubit interactions.

\section{Limitations and Threats to Validity}

The experiment is a simulation study on one dataset and one calibrated fake
backend; it does not execute circuits on live hardware.  Backend
calibrations represent a snapshot and results may change across devices or
dates.  Only one local observable and the first ansatz parameter are used
for the gradient statistic, so the metric does not characterize the full
optimization landscape.  Ten seeds bind ten different training examples,
confounding initialization variability with input variability; future work
should cross data samples and initialization seeds independently.

No end-to-end classifier is trained, so ``trainability'' here means gradient
survival rather than achieved accuracy, convergence time, or generalization.
The noisy Monte Carlo budget is smaller than the noiseless budget, which
increases uncertainty in ratios near zero.  Exact density-matrix simulation
also limits the study to ten qubits.  Finally, one of 192 scheduled
configurations is missing and is excluded without imputation.  Conclusions
about full-topology, ten-qubit, depth-eight behavior should therefore be
read with that missing \texttt{local\_z}+\texttt{real\_amplitudes} cell in
mind.

\section{Conclusion}

This crossed experiment shows that expressibility and trainability cannot be
assigned to a feature map or ansatz in isolation.  Under noiseless
simulation, increasing depth strongly suppresses gradients for the
entangling \texttt{real\_amplitudes} ansatz, while a rotation-only ansatz can
remain trainable even at depth eight.  Under a fixed calibrated noise model,
logical topology becomes decisive: linear circuits retain useful parameter
sensitivity, whereas routing circular and fully connected circuits onto a
sparse patch often produces near-complete gradient collapse.  The practical
conclusion is a hardware-aware one: connectivity-compatible shallow circuits
should be screened before paying for full variational training, and the
screen must evaluate the composed encoding--ansatz circuit rather than a
single component.

\bibliographystyle{IEEEtran}
\bibliography{references}

\end{document}
